\chapter{Conclusion \& Future Outlook}
\label{ch:conclusion}

%==============================================================================
\section{Concluding Remarks}
\label{sec:conclusion}

\draftnote{Written last, once Chapter~\ref{ch:results} is filled and the headline framing is
  settled. The skeleton below records what the evidence supports as of drafting, so that the
  final version can be checked against it rather than drifting.}

This thesis set out to predict conversion from mild cognitive impairment to Alzheimer's
dementia from longitudinal resting-state functional connectomes, using a staged pipeline of
self-supervised graph pretraining followed by a time-conditioned recurrent trajectory
classifier. What it delivers is partly that, and partly a set of findings about the
methodology of the field it sits in.

Three results are robust to the caveats stated throughout.

\begin{enumerate}
  \item \textbf{A published state-of-the-art baseline is not reproducible at a fixed seed},
    and the standard multi-seed error bar misattributes its variance. Across 19 runs of one
    configuration, within-seed standard deviation ($0.1011$) exceeds between-seed standard
    deviation ($0.0706$), and one seed spans an AUC range of $0.35$. The pipeline developed
    here, by contrast, re-runs byte-identically.
  \item \textbf{The graph encoder does not earn its place on this cohort.} Mean-pooling the
    raw connectivity features and feeding them directly to the recurrent core matches the
    frozen pretrained encoder at roughly seventeen times fewer trainable parameters ---
    extending a published cross-sectional, large-$N$ result into the longitudinal, small-$N$
    regime where it had not been tested.
  \item \textbf{Time-conditioning cannot be evidenced on a protocol-driven cohort.} Ninety
    per cent of DELCODE's inter-visit intervals are identical, and ablating the $\Delta t$
    input changes held-out AUC by exactly zero. Visit \emph{order} does carry signal;
    inter-visit \emph{timing}, on this cohort, cannot.
\end{enumerate}

Taken together these point somewhere slightly different from where the project began. The
useful contribution is less ``a better architecture for brain graphs'' and more ``an
evaluation protocol under which the architecture question can be answered honestly, plus the
answers it gives on this cohort.''

%==============================================================================
\section{Future Research Directions}
\label{sec:future}

\paragraph{Complete the external replication.} ADNI ($162$ longitudinal subjects, $51$
converters) and OASIS-3 ($60$ / $31$) are preprocessed and waiting. Both encode genuine
elapsed time rather than protocol months, which makes them the only place the
time-conditioning question of Section~\ref{sec:delta-t-finding} can be answered. This is the
highest-value open item by a wide margin.

\paragraph{Test the graph construction, not just the encoder.} The negative encoder result
was obtained under $k$NN-8 sparsification with binary edge attributes --- a construction that
discards correlation magnitude before message passing
(Section~\ref{sec:graph-construction}). A weighted-edge or fully-dense variant would
separate ``message passing does not help'' from ``this graph does not carry the
information''. That is a cheap and directly informative experiment.

\paragraph{Replace the pretext objective.} If reconstruction supplies no gradient preserving
diagnostically relevant between-subject variance (Section~\ref{sec:gap-pretext}), the
principled response is a contrastive or joint-embedding predictive objective rather than a
larger reconstruction model. Whether such objectives work at $N < 150$, away from the
foundation-model scale where they were developed, is open.

\paragraph{The parcellation ablation.} Whole brain versus DMN versus DMN with limbic and
hippocampal regions was planned and not run (Section~\ref{sec:res-not-run}). Given that DMN
dysconnectivity is the most replicated functional finding in early AD, a region-restricted
model is a natural test of whether the pipeline is using disease-relevant signal.

\paragraph{The time-to-event head.} Kaplan--Meier and clinical Cox baselines exist; an
imaging-based survival head consuming the same latents does not
(Section~\ref{sec:prognoser-scope}). This is what converts a binary flag into the calibrated
horizon risk a clinician can act on.

\paragraph{Multimodal fusion.} Plasma p-tau217 has changed what is cheaply measurable in
prodromal AD since this project began. Fusing a blood biomarker, structural T1, and
functional connectivity is the direction with the clearest route to clinical relevance ---
and, given the results of Chapter~\ref{ch:results}, the honest first question to ask of it
is whether the imaging branch adds anything over the blood test alone.

\paragraph{Adopt same-seed replication as protocol.} The cheapest and most transferable
recommendation from this work: report within-seed variance alongside between-seed variance,
always. It costs a few runs and would have prevented an invalid significance claim in this
project, in a published baseline, and --- if the argument of
Section~\ref{sec:gap-determinism} is right --- in a substantial fraction of the brain-graph
literature.
